\chapter{Introduction}
Predicting molecular properties directly from chemical structure \textit{in silico} is a central objective in computational chemistry, with applications spanning materials discovery, drug development, and molecular design. In recent years, machine learning methods operating on molecular graphs or Simplified Molecular Input Line Entry System (SMILES) strings have achieved strong empirical performance. By learning representations directly from molecular structure, these models aim to replace or augment expensive quantum-chemical simulations with scalable, data-driven surrogates.

Despite these advances, many chemically relevant properties remain challenging to predict reliably. Optical absorption and emission, in particular, depend not only on dominant electronic factors such as energy gaps and $\pi$-conjugation, but also on subtler structural and environmental effects. Datasets for these properties are often limited in size, heterogeneous in composition, and influenced by experimental variability. This inherent data scarcity limits the statistical efficiency of high-capacity neural models and increases sensitivity to training stochasticity.

At a structural level, molecules are naturally represented as graphs, where atoms correspond to nodes and chemical bonds correspond to edges. Graph-based representations preserve topological relationships while maintaining invariance to atom indexing and molecular orientation. Graph neural networks (GNNs) operationalize this structure by propagating information along bonds through iterative message-passing layers, constructing latent embeddings that summarize local chemical environments and their interactions. 

Most molecular machine learning pipelines rely on end-to-end architectures that directly couple these learned embeddings with a neural ``head'' that maps those embeddings to the predicted response. While effective in data-rich domains and required for jointly training an encoder-head pair via backpropagation, this paradigm implicitly merges representation learning and prediction into a single optimization problem. This coupling makes it difficult to determine whether improvements arise from better structural embeddings, from flexible downstream regression, or from favourable stochastic variation. While theoretically, this jointly optimized encoder-head pair can perfectly fit the data, in practical low-data regimes, the high capacity of these models can lead to overfitting, with regularization becoming nontrivial. 

These vulnerabilities become even more evident when there is a distribution shift. Scaffold-based splits, which group molecules by their core chemical structure, provide a better approximation of real-world generalization to new chemical spaces. However, they often show a significant drop in performance compared to random splits. In such cases, relying on average performance metrics can mask the considerable variability observed across different training runs.

Due to the current practical limitations of deep learning, alternative modelling paradigms remain widely used. Fixed molecular descriptors paired with classical regression methods, as well as physical approximations like time-dependent density functional theory (TD-DFT), offer competitive statistical efficiency and strong inductive biases in data-scarce environments. The coexistence of these approaches underscores a need to understand not only which models perform well, but under what specific structural constraints expressive neural modelling actually provides tangible benefits over modular alternatives.

Motivated by these challenges, this thesis explores a representation-centric approach to optical property prediction. Rather than focusing solely on end-to-end predictive accuracy, this work examines the behaviour of molecular representations learned by GNNs, particularly in data-scarce environments. By decoupling the graph encoder from downstream regression tasks and pairing it with implicitly regularized tree-based models, we investigate hybrid modelling strategies for improving performance in data-limited regimes.

The contributions of this thesis are fourfold. First, it systematically compares a broad range of modelling paradigms, including classical descriptor-based models, tree-based methods, physics-informed approaches, and end-to-end graph neural networks, using a Directed Message Passing Neural Network (D-MPNN) as a central reference point. The comparison spans multiple datasets, split strategies, and training set sizes. Second, it evaluates hybrid modelling strategies, demonstrating that decoupling representation learning from downstream prediction provides distinct advantages in low-data regimes. Third, it systematically analyzes the sources of predictive variability using nested variance decomposition, providing a mathematically grounded view of instability. Finally, it investigates the learned latent representations to assess whether the graph encoders capture chemically meaningful structure--property relationships or merely exploit dataset-specific artifacts. These contributions are primarily analytical rather than architectural. Rather than proposing a completely novel model, this work provides a practical characterization of when different modelling approaches are effective for molecular optical property prediction in data-limited settings typical of experimental labs, and identifies the sources of their performance differences in terms of representation quality, downstream regression, and stochastic variability.

The remainder of this thesis proceeds as follows. Chapter~\ref{chap:related_work} situates the present work within the literature on molecular graph learning, molecular property prediction, transfer learning, and evaluation under data scarcity. Chapter~\ref{chap:math_foundations} formalizes the mathematical framework underlying representation learning with MPNNs, our choice of tree-based regressor, and variance decomposition. Chapter~\ref{chap:manuscript} presents the experimental methodology and empirical findings, with additional supporting analysis provided in Chapter~\ref{chap:supporting}. Chapter~\ref{chap:conclusion} synthesizes the results and discusses broader implications.
